%! Unit 7 — Analytic Trigonometry and Polar Coordinates — Unit Cover Page
\documentclass[10pt]{article}
\usepackage{precalculus-article}
\usepackage{precalculus-boxes}
\usepackage{ltablex}
\keepXColumns

\begin{document}
\thispagestyle{empty}

% Full-bleed navy banner
\begin{tikzpicture}[remember picture, overlay]
  \fill[navy] (current page.north west)
    rectangle ([yshift=-1.15in]current page.north east);
\end{tikzpicture}
\vspace{-1.05in}

\begin{center}
  {\color{white}\textbf{\LARGE Precalculus \quad 2026--2027}}\\[6pt]
  {\color{skymid}\Large\textbf{Unit 7: Analytic Trigonometry and Polar Coordinates}}\\[4pt]
  {\color{white}\small 8 Lessons + Unit Test \quad|\quad Class Periods: 18--22}
\end{center}

\vspace{0.26in}

\begin{tcolorbox}[colback=goldbg, colframe=goldacc, boxrule=0.8pt, arc=2mm,
    left=3mm, right=3mm, top=2mm, bottom=2mm]
\small\textbf{Unit Essential Questions:}
\begin{itemize}[leftmargin=1.4em, itemsep=2pt, topsep=2pt]
  \item How do inverse trigonometric functions and identities let us solve trigonometric equations exactly?
  \item What does the polar coordinate system reveal that Cartesian coordinates hide?
\end{itemize}
\end{tcolorbox}

\vspace{0.14in}

\begin{tocbox}
\small
\renewcommand{\arraystretch}{1.45}
\begin{tabularx}{\linewidth}{@{} c >{\bfseries}X l @{}}
\rowcolor{sky}
\textbf{\#} & \textbf{Lesson Title} & \textbf{Content Source} \\
\midrule
7.1 & The Tangent Function & CED 3.8 \\
7.2 & Inverse Trigonometric Functions & CED 3.9 \\
7.3 & Trigonometric Equations and Inequalities & CED 3.10 \\
7.4 & The Secant, Cosecant, and Cotangent Functions & CED 3.11 \\
7.5 & Equivalent Representations of Trigonometric Functions & CED 3.12 \\
7.6 & Trigonometry and Polar Coordinates & CED 3.13 \\
7.7 & Polar Function Graphs & CED 3.14 \\
7.8 & Rates of Change in Polar Functions & CED 3.15 \\
\midrule
\rowcolor{sky}
     & \textbf{Sample Test} &      \\
\end{tabularx}
\end{tocbox}

\vspace{0.14in}

\begin{skillbox}[Mathematical Practices --- Unit 7 Focus]{goldbox}
\renewcommand{\arraystretch}{1.4}
\begin{tabularx}{\linewidth}{@{} >{\bfseries\color{navy}}p{0.12\linewidth} X @{}}
MP 1.A & Solve equations and inequalities analytically, with and without technology. \\
MP 1.B & Express functions, equations, or expressions in analytically equivalent forms. \\
MP 2.B & Construct equivalent graphical, numerical, analytical, and verbal representations. \\
MP 3.A & Describe the characteristics of a function with varying levels of precision. \\
\end{tabularx}
\end{skillbox}

\end{document}
